% ==============================================================
%  Chapter 1: The AI Landscape
% ==============================================================
\chapter{The AI Landscape: What Exists and How to Use It}
\label{ch:overview}

The AI ecosystem has expanded faster than any technology shift since the smartphone.
There are now hundreds of tools, but they fall into a small number of clear categories.
Understanding \emph{which tool does what} is the first skill.
Using the wrong tool --- asking a chat model for real-time stock prices, or an image
generator to review code --- wastes time and produces bad results.
This chapter is your map.

\section{Categories of AI Tools}
\label{sec:categories}

\begin{table}[htbp]
\centering
\caption{AI tool categories, primary strengths, and representative examples}
\begin{tabularx}{\linewidth}{l l X}
\toprule
\textbf{Category} & \textbf{Examples} & \textbf{Best For} \\
\midrule
Conversational LLMs    & ChatGPT, Claude, Gemini   & Reasoning, writing, Q\&A, planning, summarisation \\[4pt]
Code Assistants        & GitHub Copilot, Cursor, Cody & Inline completion, refactoring, codebase chat \\[4pt]
Agentic AI             & Copilot Agent, Cursor Agent, Devin & Multi-step autonomous task execution \\[4pt]
Search AI              & Perplexity, Bing Copilot  & Research with live citations \\[4pt]
Image / Video AI       & Midjourney, DALL-E, Runway & Diagrams, mockups, visual content \\[4pt]
Voice / Transcript AI  & Whisper, Otter.ai, Teams Copilot & Meeting notes, audio transcription \\[4pt]
Data Analysis AI       & Code Interpreter, Julius  & EDA, ad-hoc analysis, chart generation \\[4pt]
RAG / Knowledge AI     & Notion AI, Confluence AI, custom & Q\&A over your own documents \\[4pt]
Workflow Agents        & CrewAI, LangGraph, n8n    & Chained AI steps, tool use, pipelines \\
\bottomrule
\end{tabularx}
\label{tab:ai-categories}
\end{table}

Each category has a different strength \emph{and} a different failure mode.
Chat LLMs reason and write well but hallucinate facts.
Search AI cites real sources but is shallower on complex reasoning.
Code assistants stay in your flow but cannot see the big picture.
Knowing the failure mode of your tool prevents expensive surprises.

\section{One-Sentence Tool Selection Guide}
\label{sec:choosing}

\begin{itemize}
  \item \textbf{Explain, reason, plan, write a draft} $\to$ conversational LLM (Claude, ChatGPT)
  \item \textbf{Inline code while you type} $\to$ code assistant (Copilot, Cursor)
  \item \textbf{Current facts with verifiable sources} $\to$ search AI (Perplexity)
  \item \textbf{Your company's internal documents} $\to$ RAG tool (Notion AI, Confluence AI)
  \item \textbf{Generate an image or diagram} $\to$ Midjourney or DALL-E
  \item \textbf{Transcribe or summarise a meeting} $\to$ Otter.ai or Teams Copilot
  \item \textbf{Execute a multi-step automated task} $\to$ agent mode (Cursor Agent, n8n)
  \item \textbf{Explore and visualise a CSV} $\to$ ChatGPT Code Interpreter or Julius
\end{itemize}

\begin{tipbox}
Start with \textbf{one tool per use case} rather than evaluating ten simultaneously.
Deep proficiency with two tools produces more value than shallow use of ten.
\end{tipbox}

\section{The Capability Spectrum}
\label{sec:spectrum}

AI capability is not uniform. Think of it as a spectrum:

\textbf{Strong:} Pattern recognition in well-documented domains (code, writing, common
knowledge), translation between formats, summarisation, first-draft generation, explaining
concepts with analogies, refactoring existing structures.

\textbf{Moderate:} Multi-step reasoning, novel combinations of ideas, identifying
unstated assumptions, evaluating tradeoffs, generating options in complex design
problems.

\textbf{Weak:} Highly novel problems with no prior art, strict factual accuracy under
time pressure, knowing what it doesn't know, reasoning about real-world state it hasn't
been shown, understanding implicit organisational or social context.

The better you understand this spectrum, the better you can direct AI toward what it
does well and take over where it struggles.

\section{How to Maximise Value from AI}
\label{sec:maximise}

The professionals who extract the most value from AI share three habits:

\begin{enumerate}
  \item \textbf{Clear context before asking.}
    They tell the AI what role it plays, the constraints, the goal, and the format they
    want --- before asking anything. An AI that knows it is acting as a senior DevOps
    engineer reviewing infrastructure security gives a different (better) answer than one
    given a vague question.

  \item \textbf{Iterative refinement, not one-shot hoping.}
    They treat the first output as a working draft and refine with follow-up prompts
    (\textit{``make it shorter'', ``add an example'', ``explain why this is better than X''}).
    One-shot expectations lead to disappointment; iteration leads to quality.

  \item \textbf{Critical evaluation before use.}
    They read the output before using it. They run code, check facts, question logic,
    and test edge cases. AI output is a fast first draft, not a final deliverable.
\end{enumerate}

\section{What AI Fundamentally Cannot Do}
\label{sec:limitations}

Understanding limitations is more valuable than understanding capabilities, because
capabilities are obvious while limitations are where costly mistakes happen.

\begin{warnbox}[Hard Limitations]
\begin{itemize}
  \item Cannot access data that appeared after its training cutoff
  \item Cannot access real-time information without explicit tool integration
  \item Cannot know your specific codebase, business rules, or org context unless you
        provide them
  \item Cannot guarantee factual accuracy --- it generates plausible text, not verified facts
  \item Cannot reason about highly novel, first-principles problems without guidance
  \item Cannot replace domain expertise, professional judgment, or accountability
  \item Will agree with you if you push back, even when it was originally correct
\end{itemize}
\end{warnbox}

AI is a force multiplier for people who already have skills, not a substitute for
developing them. The more domain knowledge you bring, the better you direct AI and
catch its errors.

\section{The Human--AI Collaboration Model}
\label{sec:collaboration}

Think of AI as a very fast, very well-read junior colleague who:

\begin{itemize}
  \item Has read almost everything published but has experienced nothing directly
  \item Can produce a confident first draft of almost anything in seconds
  \item Sometimes makes up facts with complete conviction
  \item Has no stake in the outcome and will tend to agree with whoever is talking
  \item Works best with clear specifications, constraints, and feedback loops
\end{itemize}

Your job as the human is to be the senior: set direction, supply context, evaluate
output critically, and take full responsibility for what gets shipped.
